\section{Controlled Demonstrations}\label{sec:results}

We now test the account in three settings where the ground truth is known. The Markov benchmark checks the exact theory and the route-mismatch diagnostic. The budget benchmark measures how much of a known deficit packaged memory can recover. The hashing experiment carries the same logic to a deterministic map. All numbers below come from the tracked outputs in the companion repository and from an exact-enumeration script that computes the population quantities.

\subsection{Markov coarse-graining benchmark}\label{sec:results-markov}

\paragraph{Design.} The benchmark contains four families of finite Markov chains, each with a known packaging into three packages.
\begin{itemize}
\item \emph{Exactly lumpable} chains (six microstates, three packages of two) are closed by construction.
\item \emph{Perturbed-lumpable} chains start from the same construction. A heterogeneity parameter $\alpha\in\{0,0.2,0.5\}$ then mixes a microstate-specific random row into each microstate's package-level transition probabilities.
\item \emph{Metastable} chains (eight microstates in blocks of three, three, and two) move mostly within a block. An escape-weight parameter $p_{\mathrm{out}}\in\{0.02,0.08,0.20\}$ controls the probability of leaving it.
\item \emph{Hidden-type} chains (eight microstates in packages of three, two, and three) split each package into two microstate types. A strength parameter $s\in\{0.1,0.5,0.9\}$ sets how differently the two types move between packages.
\end{itemize}
Every kernel is mixed with a small uniform component, so all transition probabilities are positive and the stationary law is unique with full support. Lumpable chains stay lumpable under this mixing. We evaluate every family at lags $\tau=1$ and $\tau=2$, except hidden types, which we evaluate at $\tau=1$ only. This gives seventeen conditions. For each one, we compute the deficit, the intrinsic term, the residual of the decomposition \eqref{eq:cd-decomposition}, and both route mismatches by floating-point matrix calculations from the known kernels.

\paragraph{Closed cases and numerical checks.} The exactly lumpable chains give deficits of magnitude below $7.6\times10^{-17}$ and uniform-lift mismatches below $2.5\times10^{-17}$ at both lags, which is floating-point zero, as \Cref{prop:cd-zero} requires. The decomposition holds to within $1.7\times10^{-16}$ in every condition, and the Pinsker bound of \Cref{prop:rm-lower-bound} holds in every condition up to roundoff of the same size. These checks confirm the implementation; they do not replace the proofs.

\paragraph{Open cases.} Once the construction breaks closure, the deficit is positive and grows with the control parameter (\Cref{fig:markov-knobs}). In the perturbed-lumpable family at $\tau=1$, raising $\alpha$ from $0.2$ to $0.5$ raises the deficit from $0.0083$ to $0.048$ nats. The strongest hidden-type chain reaches $0.031$ nats. The metastable family stays small, rising to $7.2\times10^{-4}$ nats at $p_{\mathrm{out}}=0.20$. There, each microstate's escape weight carries its own random multiplier, so the spread of exit probabilities within a block grows with $p_{\mathrm{out}}$.

\begin{figure}[tbp]
\centering
\includegraphics[width=\linewidth]{fig_markov_knobs}
\caption{Closure deficit against each family's control parameter, at lag $\tau=1$ (filled markers, solid lines) and $\tau=2$ (open markers, dashed lines). Each panel has its own vertical scale. The deficit is zero when the control parameter removes the hidden distinction ($\alpha=0$) and grows as the distinction strengthens. For the open conditions evaluated at both lags, the lag-two deficit is smaller than the lag-one deficit, but that is not a general law (\Cref{rem:nonmonotone}).}
\label{fig:markov-knobs}
\end{figure}

\paragraph{Route mismatch as a proxy.} Across all seventeen conditions, the uniform-lift mismatch has Pearson correlation $r=0.959$ with the deficit, and the stationary-lift mismatch has $r=0.961$. \Cref{fig:markov-rm-cd} shows why on logarithmic axes. Across more than three orders of magnitude in deficit, both mismatches rank the thirteen conditions with nonzero deficit nearly as the deficit does (Spearman rank correlation $0.984$ for the uniform lift and $0.989$ for the stationary lift), and every point lies above the Pinsker curve $\tfrac12(\mathrm{RM}^{\mathrm{stat}})^2$. In the most heterogeneous perturbed chain, for example, $\mathsf{CD}_1=0.048$ with $\mathrm{RM}^{\mathrm{unif}}_1=0.276$ and $\mathrm{RM}^{\mathrm{stat}}_1=0.270$. For the uniform lift, this agreement is empirical. \Cref{rem:uniform-lift} shows that it can fail when a fiber contains a rarely visited microstate. In this exact implementation, both mismatches use the same destination rows and stationary law as the deficit itself, so the benchmark establishes an association, not a computational shortcut.

\begin{figure}[tbp]
\centering
\includegraphics[width=\linewidth]{fig_markov_rm_cd}
\caption{Closure deficit against route mismatch for the thirteen conditions with nonzero deficit, on logarithmic axes. The four numerically closed conditions (exactly lumpable chains and $\alpha=0$, at both lags) are at floating-point zero and are omitted. (a)~Stationary-lift mismatch; the curve $\tfrac12(\mathrm{RM}^{\mathrm{stat}})^2$ is the guaranteed lower bound of \Cref{prop:rm-lower-bound}. (b)~Uniform-lift mismatch, for which no such bound holds in general. Filled markers are lag $\tau=1$ and open markers lag $\tau=2$.}
\label{fig:markov-rm-cd}
\end{figure}

\paragraph{Lag.} In both open families evaluated at two lags, the deficit at $\tau=2$ is smaller than at $\tau=1$ for every open condition ($\alpha>0$ and every $p_{\mathrm{out}}$). For $\alpha=0.2$ and $0.5$, it falls from $0.0083$ to $0.0012$ and from $0.048$ to $0.0048$. In these strongly mixing chains, two steps of dynamics wash out much of the hidden distinction before the next package is recorded. This is a property of these chains, not of staging in general. The four-cycle of \Cref{rem:nonmonotone} has zero deficit at lag two and $\log 2$ at lags one and three. What the benchmark supports is the Six Birds point about staging (P4): randomness at a layer depends on the time scale at which the layer is run, and the dependence must be computed, not assumed.
